\section{Khoảng biến thiên và khoảng tứ phân vị của mẫu số liệu ghép nhóm}
\subsection{Đề 1}
\Opensolutionfile{ans}[ans/ans-0-B15]
\caulc
\begin{ex}%[2D3N1-1]
 Phát biểu nào sau đây là \textbf{sai}
\choice{Khoảng tứ phân vị càng lớn thì mẫu số liệu càng phân tán}
{Khoảng tứ phân vị không phụ thuộc vào các giá trị bất thường}
{Khoảng biến thiên càng bé thì độ phân tán càng bé}
{\True Khoảng biến thiên không phụ thuộc vào các giá trị bất thường}
\loigiai{
Nếu mẫu số liệu chứa giá trị bất thường, quá cao hoặc quá thấp thì khoảng biến thiên sẽ thay đổi nhiều. Do đó khoảng biến thiên phụ thuộc vào các giá trị bất thường.
 }
\end{ex}
\begin{ex}%[1D5N1-2]
Cho mẫu số liệu ghép nhóm cho bởi bảng như hình sau
\begin{center}
	\begin{tabular}{|c|c|c|c|c|c|}
		\hline
		\textbf{Nhóm} & $[160;163)$&$[163;166)$& $[166;169)$& $[169;172)$&$[172;175)$\\
		\hline
		\textbf{Tần số} & $6$ &$11$& $9$& $7$&$3$\\
		\hline
	\end{tabular}
\end{center}
Cỡ mẫu của mẫu số liệu bằng
\choice
{$11$}
{\True $36$}
{$175$}
{$15$}
\loigiai{
Cỡ mẫu của mẫu số liệu là $n=6+11+9+7+3=36$.
}
\end{ex}
\begin{ex}%[1D5N1-2]
Cho mẫu số liệu ghép nhóm cho bởi bảng như hình sau
	\begin{center}
		\begin{tabular}{|c|c|c|c|c|c|}
			\hline
			\textbf{Nhóm} & $[160;163)$&$[163;166)$& $[166;169)$& $[169;172)$&$[172;175)$\\
			\hline
			\textbf{Tần số} & $6$ &$11$& $9$& $7$&$3$\\
			\hline
		\end{tabular}
	\end{center}
Độ dài mỗi nhóm bằng
\choice
{$2$}
{\True $3$}
{$5$}
{$4$}
\loigiai{Độ dài mỗi nhóm bằng $\ell =3$.}
\end{ex}
\begin{ex}%[2D3H1-2]
Cho mẫu số liệu ghép nhóm cho bởi bảng như hình sau
\begin{center}
	\begin{tabular}{|c|c|c|c|c|c|}
		\hline
		\textbf{Nhóm} & $[160;163)$&$[163;166)$& $[166;169)$& $[169;172)$&$[172;175)$\\
		\hline
		\textbf{Tần số} & $6$ &$11$& $9$& $7$&$3$\\
		\hline
	\end{tabular}
\end{center}
Khoảng biến thiên của mẫu số liệu là
\choice{$3$}{$9$}{$8$}{\True $15$}
\loigiai{
Khoảng biến thiên của mẫu số liệu là $R=175-160=15$.
}
\end{ex}
\begin{ex}%[1D5N1-4]
Cho mẫu số liệu ghép nhóm cho bởi bảng như hình sau
	\begin{center}
		\begin{tabular}{|c|c|c|c|c|c|}
			\hline
			\textbf{Nhóm} & $[160;163)$&$[163;166)$& $[166;169)$& $[169;172)$&$[172;175)$\\
			\hline
			\textbf{Tần số} & $6$ &$11$& $9$& $7$&$3$\\
			\hline
		\end{tabular}
	\end{center}
Khoảng chứa mốt của mẫu số liệu là
	\choice
	{$[163;166)$}
	{\True $[166;169)$}
	{$[169;172)$}
	{$[172;175)$}
	\loigiai{		
	}
\end{ex}
\begin{ex}%[1D5H1-4]
Cho mẫu số liệu ghép nhóm cho bởi bảng như hình sau
\begin{center}
	\begin{tabular}{|c|c|c|c|c|c|}
		\hline
		\textbf{Nhóm} & $[160;163)$&$[163;166)$& $[166;169)$& $[169;172)$&$[172;175)$\\
		\hline
		\textbf{Tần số} & $6$ &$11$& $9$& $7$&$3$\\
		\hline
	\end{tabular}
\end{center}
Mốt của mẫu số liệu là
	\choice
	{$166$}
	{$163$}
	{\True $165{,}14$}
	{$166{,}6$}
	\loigiai{
Mốt của mẫu số liệu là $M_o=163+\dfrac{5}{5+2}\times 3=165{,}14$.		
	}
\end{ex}
\begin{ex}%[1D5H2-3]
Cho mẫu số liệu ghép nhóm cho bởi bảng như hình sau
\begin{center}
	\begin{tabular}{|c|c|c|c|c|c|}
		\hline
		\textbf{Nhóm} & $[160;163)$&$[163;166)$& $[166;169)$& $[169;172)$&$[172;175)$\\
		\hline
		\textbf{Tần số} & $6$ &$11$& $9$& $7$&$3$\\
		\hline
	\end{tabular}
\end{center}
Tứ phân vị thứ nhất của mẫu là
	\choice
	{$164{,}1$}
	{$163{,}2$}
	{$163{,}5$}
	{\True $163{,}8$}
	\loigiai{
Ta có $Q_1=163+\dfrac{9-5}{11}\times 3=163{,}8$.
	}
\end{ex}
\begin{ex}%[2D3V1-3]
	Cho mẫu số liệu ghép nhóm cho bởi bảng như hình sau
	\begin{center}
		\begin{tabular}{|c|c|c|c|c|c|}
			\hline
			\textbf{Nhóm} & $[160;163)$&$[163;166)$& $[166;169)$& $[169;172)$&$[172;175)$\\
			\hline
			\textbf{Tần số} & $6$ &$11$& $9$& $7$&$3$\\
			\hline
		\end{tabular}
	\end{center}
Khoảng tứ phân vị của mẫu số liệu là
	\choice
	{\True $5{,}6$}
	{$5{,}2$}
	{$6{,}4$}
	{$6{,}8$}
	\loigiai{
Ta có $Q_1=163+\dfrac{9-5}{11}\times 3=163{,}8$.\\
Ta có $Q_3=169+\dfrac{27-26}{7}\times 3=169{,}4$. Do đó khoảng tứ phân vị là $\Delta Q=Q_3-Q_1=5{,}6$.			
	}
\end{ex}
\begin{ex}%[1D5V1-3]
Biểu đồ dưới đây thống kê thời gian tập thể dục buổi sáng mỗi ngày trong tháng 9/2022 của bác Bình và bác An.
\begin{center}
	\begin{tikzpicture}[xscale=2,font=\footnotesize, line join=round, line cap=round, >=stealth]
		\draw[->](0,-0.2)--(0,7)node[left]{\textbf{Số ngày}};
		\draw[->](-0.2,0)--(6,0)node[below]{\textbf{phút}};
		\foreach \i in{1,...,6} \pgfmathsetmacro{\gti}{int(5*(\i))}
		\draw [dotted](0,\i) circle(1pt)node[left]{$\gti$} -- (5.5,\i);
		\foreach \i/\a/\b in {1/15/20,2/20/25,3/25/30,4/30/35,5/35/40}\draw (\i,0) node [below] {[\a;\b)};
		\foreach \i/\tr/\ph in {1/5/0,2/12/25,3/8/5,4/3/0,5/2/0}
		{
			\draw[fill=blue](\i,0)rectangle(\i-0.5,0.2*\tr) node [above right=0.3]{\tr};
			\draw[fill=orange](\i,0)rectangle(\i+0.5,0.2*\ph) node [above left=0.3]{\ph};
		}
		\foreach \i/\a/\b in {3/orange/Bác An, 4/blue/Bác Bình}
		{\draw (6,\i+1)node [right,fill=\a] {};
			\draw (6.2,\i+1) node[right]{\b};
		}
	\end{tikzpicture}
\end{center}
Trung bình thời gian tập thể dục mỗi ngày của bác Bình là
\choice
{$26 \text{ phút}$}
{$25{,}5 \text{ phút}$}
{$25{,}75 \text{ phút}$}
{\True $25 \text{ phút}$}
\loigiai{
Dựa vào biểu đồ ta có bảng tần số thời gian tập thể dục của bác Bình là
\begin{center}
	\begin{tabular}{|c|c|c|c|c|c|}
		\hline
		\textbf{Nhóm} & $[15;20)$&$[20;25)$& $[25,30)$& $[30;35)$&$[35;40)$\\
		\hline
		\textbf{Tần số} & $5$ &$12$& $8$& $3$&$2$\\
		\hline
		\textbf{Giá trị đại diện} & $17{,}5$ &$22{,}5$& $27{,}5$& $32{,}5$&$37{,5}$\\
		\hline		
	\end{tabular}
\end{center}
Dựa vào bảng trên ta có giá trị trung bình bằng $$\dfrac{17{,}5\times5+22{,}5\times12+27{,}5\times8+32{,}5\times3+37{,}5\times2}{30}=25 \text{ phút}.$$		
	}
\end{ex}
\begin{ex}%[2D3H1-2]
Biểu đồ dưới đây thống kê thời gian tập thể dục buổi sáng mỗi ngày trong tháng 9/2022 của bác Bình và bác An.
\begin{center}
	\begin{tikzpicture}[xscale=2,font=\footnotesize, line join=round, line cap=round, >=stealth]
		\draw[->](0,-0.2)--(0,7)node[left]{\textbf{Số ngày}};
		\draw[->](-0.2,0)--(6,0)node[below]{\textbf{phút}};
		\foreach \i in{1,...,6} \pgfmathsetmacro{\gti}{int(5*(\i))}
		\draw [dotted](0,\i) circle(1pt)node[left]{$\gti$} -- (5.5,\i);
		\foreach \i/\a/\b in {1/15/20,2/20/25,3/25/30,4/30/35,5/35/40}\draw (\i,0) node [below] {[\a;\b)};
		\foreach \i/\tr/\ph in {1/5/0,2/12/25,3/8/5,4/3/0,5/2/0}
		{
			\draw[fill=blue](\i,0)rectangle(\i-0.5,0.2*\tr) node [above right=0.3]{\tr};
			\draw[fill=orange](\i,0)rectangle(\i+0.5,0.2*\ph) node [above left=0.3]{\ph};
		}
		\foreach \i/\a/\b in {3/orange/Bác An, 4/blue/Bác Bình}
		{\draw (6,\i+1)node [right,fill=\a] {};
			\draw (6.2,\i+1) node[right]{\b};
		}
	\end{tikzpicture}
\end{center}
Khoảng biến thiên biểu thị thời gian tập thể dục của bác An là
	\choice
	{\True $10$ phút}
	{$15$ phút}
	{$20$ phút}
	{$25$ phút}
	\loigiai{
Khoảng biến thiên của bác An là $R=30-20=10$ phút.
	}
\end{ex}
\begin{ex}%[1D5H2-3]
Biểu đồ dưới đây thống kê thời gian tập thể dục buổi sáng mỗi ngày trong tháng 9/2022 của bác Bình và bác An.
\begin{center}
	\begin{tikzpicture}[xscale=2,font=\footnotesize, line join=round, line cap=round, >=stealth]
		\draw[->](0,-0.2)--(0,7)node[left]{\textbf{Số ngày}};
		\draw[->](-0.2,0)--(6,0)node[below]{\textbf{phút}};
		\foreach \i in{1,...,6} \pgfmathsetmacro{\gti}{int(5*(\i))}
		\draw [dotted](0,\i) circle(1pt)node[left]{$\gti$} -- (5.5,\i);
		\foreach \i/\a/\b in {1/15/20,2/20/25,3/25/30,4/30/35,5/35/40}\draw (\i,0) node [below] {[\a;\b)};
		\foreach \i/\tr/\ph in {1/5/0,2/12/25,3/8/5,4/3/0,5/2/0}
		{
			\draw[fill=blue](\i,0)rectangle(\i-0.5,0.2*\tr) node [above right=0.3]{\tr};
			\draw[fill=orange](\i,0)rectangle(\i+0.5,0.2*\ph) node [above left=0.3]{\ph};
		}
		\foreach \i/\a/\b in {3/orange/Bác An, 4/blue/Bác Bình}
		{\draw (6,\i+1)node [right,fill=\a] {};
			\draw (6.2,\i+1) node[right]{\b};
		}
	\end{tikzpicture}
\end{center}
Tứ phân vị thứ nhất của mẫu số liệu thời gian tập thể dục của bác Bình là
	\choice
	{$22{,}4$}
	{$22{,}14$}
	{\True $21{,}04$}
	{$21{,}4$}
	\loigiai{
Tứ phân vị thứ nhất của mẫu số liệu thời gian tập thể dục của bác Bình là $$Q_1=20+\dfrac{7{,}5-5}{12}\times 5=21{,}04.$$		
	}
\end{ex}
\begin{ex}%[2D3V1-3]
Biểu đồ dưới đây thống kê thời gian tập thể dục buổi sáng mỗi ngày trong tháng 9/2022 của bác Bình và bác An.
\begin{center}
	\begin{tikzpicture}[xscale=2,font=\footnotesize, line join=round, line cap=round, >=stealth]
		\draw[->](0,-0.2)--(0,7)node[left]{\textbf{Số ngày}};
		\draw[->](-0.2,0)--(6,0)node[below]{\textbf{phút}};
		\foreach \i in{1,...,6} \pgfmathsetmacro{\gti}{int(5*(\i))}
		\draw [dotted](0,\i) circle(1pt)node[left]{$\gti$} -- (5.5,\i);
		\foreach \i/\a/\b in {1/15/20,2/20/25,3/25/30,4/30/35,5/35/40}\draw (\i,0) node [below] {[\a;\b)};
		\foreach \i/\tr/\ph in {1/5/0,2/12/25,3/8/5,4/3/0,5/2/0}
		{
			\draw[fill=blue](\i,0)rectangle(\i-0.5,0.2*\tr) node [above right=0.3]{\tr};
			\draw[fill=orange](\i,0)rectangle(\i+0.5,0.2*\ph) node [above left=0.3]{\ph};
		}
		\foreach \i/\a/\b in {3/orange/Bác An, 4/blue/Bác Bình}
		{\draw (6,\i+1)node [right,fill=\a] {};
			\draw (6.2,\i+1) node[right]{\b};
		}
	\end{tikzpicture}
\end{center}	

Khoảng tứ phân vị biểu thị thời gian tập thể dục của bác Bình bằng
	\choice
	{\True $7{,}4$}
	{$6{,}8$}
	{$7{,}1$}
	{$8{,}1$}
	\loigiai{
		Tứ phân vị thứ nhất của mẫu số liệu thời gian tập thể dục của bác Bình là $$Q_1=20+\dfrac{7{,}5-5}{12}\times 5=21{,}04.$$
Tứ phân vị thứ ba của mẫu số liệu thời gian tập thể dục của bác Bình là $$Q_3=25+\dfrac{22{,}5-17}{8}\times 5=28{,}44.$$
Do đó khoảng tứ phân vị của mẫu số liệu bằng $R=28{,}44-21{,}04=7{,}4$.
	}
\end{ex}
\Closesolutionfile{ans}
\indapan{6}{ans/ans-0-B15}
\cauds
\Opensolutionfile{ans}[ans/ans-0-B15-DS]
\begin{ex}%[2D3V1-2]
Bảng bên dưới biểu diễn mẫu số liệu ghép nhóm về chiều cao (đơn vị: centimét) của các học sinh nam lớp 12A và lớp 12B ở một trường trung học phổ thông.
\begin{center}
\begin{tabular}{|c|c|c|c|c|c|}
\hline
\textbf{Chiều cao} & $[160;163)$ & $[163; 166)$ & $[166;169)$ & $[169;172)$ & $[172;175)$\\
\hline
\textbf{Số học sinh nam lớp 12A} & $6$ & $11$ & $9$ & $7$ & $3$ \\
\hline
\textbf{Số học sinh nam lớp 12B} & $0$ & $21$ & $8$ & $7$ & $0$\\
\hline
\end{tabular}
\end{center}
\choiceTF
{\True Độ dài mỗi nhóm bằng $3$ (cm)}
{\True Khoảng biến thiên của mẫu số liệu ghép nhóm về chiều cao của các bạn học sinh nam lớp 12A là $15$\,(cm)}
{Khoảng biến thiên của mẫu số liệu ghép nhóm về chiều cao của các bạn học sinh nam lớp 12B là $15$\,(cm)}
{Nếu căn cứ và khoảng biến thiên thì chiều cao của các bạn nam ở hai lớp có độ phân tán như nhau}
\loigiai{
\begin{itemchoice}
\itemch Độ dài mỗi nhóm bằng $3$. Đúng.
\itemch Đúng. Vì khoảng biến thiên của mẫu số liệu ghép nhóm về chiều cao của các bạn học sinh nam lớp 12A là $175-160=15$\,(cm).
\itemch Sai. Vì khoảng biến thiên của mẫu số liệu ghép nhóm về chiều cao của các bạn học sinh nam lớp 12B là $172-163=9$\,(cm).
\itemch Sai. Vì nếu căn cứ và khoảng biến thiên thì chiều cao của các bạn nam ở lớp 12A phân tán hơn ($15 > 9$).
\end{itemchoice}
}
\end{ex}
\begin{ex}%[2D3V1-2]
Bạn Trang thống kê lại thời gian tập thể dục buổi sáng của bác Bình và bác An ở bảng sau:
\begin{center}
\begin{tabular}{|c|c|c|c|c|c|}
\hline
\textbf{Thời gian (phút)} & $[15;20)$ & $[20;25)$ & $[25;30)$ & $[30;35)$ & $[35;40)$\\
\hline
\textbf{Số ngày tập của bác Bình} & $5$ & $12$ & $8$ & $3$ & $2$ \\
\hline
\textbf{Số ngày tập của bác An} & $0$ & $25$ & $5$ & $0$ & $0$\\
\hline
\end{tabular}
\end{center}
\choiceTF
{Trong mẫu số liệu ghép nhóm về thời gian tập thể dục buổi sáng của bác An, tứ phân vị thứ nhất bằng $22{,}5$ (phút)}
{\True Khoảng biến thiên của mẫu số liệu ghép nhóm về thời gian tập của của bác An là $25$ (phút)}
{Khoảng biến thiên của mẫu số liệu ghép nhóm về thời gian tập của của bác An là $25$ (phút)}
{\True Nếu căn cứ và khoảng biến thiên thì bác Bình có thời gian tập phân tán hơn bác An}
\loigiai{
\begin{itemchoice}
\itemch Sai. Tứ phân vị thứ nhất bằng $Q_1=20+\dfrac{15-0}{25}\times 5=23$ (phút).
\itemch Đúng. Vì khoảng biến thiên của mẫu số liệu ghép nhóm về thời gian tập của của bác An là $40-15=25$ (phút).
\itemch Sai. Vì khoảng biến thiên của mẫu số liệu ghép nhóm về thời gian tập của của bác An là $30-20=10$ (phút).
\itemch Đúng. Vì khoảng biến thiên của mẫu số liệu ghép nhóm về thời gina tập của của bác Bình nhỏ hơn khoảng biến thiên của mẫu số liệu ghép nhóm về thời gina tập của của bác An.
\end{itemchoice}
}
\end{ex}
\begin{ex}%[2D3V1-3]
Một doanh nghiệp ở địa phương A muốn hướng dịch vụ của mình đến các gia đình có mức thu nhập ở tầm trung, tức là 50\% các hộ gia đình có mức thu nhập ở chính giữa so với mức thu nhập của tất cả các hộ gia đình của địa phương nên tiến hành điều tra tổng thu nhập trong năm 2022 của một số hộ gia đình trong một địa phương này. Kết quả được ghi lại trong bảng sau:
\begin{center}
\begin{tabular}{|c|c|c|c|c|c|}
\hline
\textbf{Tổng thu nhập (triệu đồng)} & $[200;250)$ & $[250;300)$ & $[300;350)$ & $[350;400)$ & $[400;450)$\\
\hline
\textbf{Số hộ gia đình} & $24$ & $62$ & $34$ & $21$ & $9$ \\
\hline
\end{tabular}
\end{center}
\choiceTF
{\True Tứ phân vị thứ nhất của mẫu số liệu ghép nhóm là $Q_1=\dfrac{16\,175}{62}$}
{Tứ phân vị thứ ba của mẫu số liệu ghép nhóm là $Q_3=\dfrac{1\,417\,125}{62}$}
{Khoảng tứ phân vị của mẫu số liệu ghép nhóm là $\Delta_Q=\dfrac{65\,475}{31}$}
{\True Doanh nghiệp cần hướng đến các gia đình có mức thu nhập trong khoảng $\left[\dfrac{16\,175}{62};\dfrac{11\,525}{34}\right)$ (triệu đồng)}
\loigiai{
Số hộ gia đình được khảo sát (cỡ mẫu) là $n=24+62+34+21+9=150$.\\
Gọi $x_1$, $x_2$, $\ldots$, $x_{150}$ là tổng thu nhập trong năm 2022 của 150 hộ gia đình được xếp theo thứ tự không giảm.\\
Ta có $x_1$, $x_2$, $\ldots$, $x_{24} \in [200;250)$; $x_{25}$, $x_{26}$, $\ldots$, $x_{86} \in [250;300)$; $x_{87}$, $x_{88}$, $\ldots$, $x_{120} \in [300;350)$; $x_{121}$, $\ldots$, $x_{141} \in [350;400)$; $x_{142}$, $\ldots$, $x_{150} \in [400;450)$.\\
Do đó, đối với dãy số liệu $x_1$, $x_2$, $\ldots$, $x_{150}$ thì
\begin{itemchoice}
\itemch Đúng. Tứ phân vị thứ nhất $Q_1$ là $x_{38} \in [250;300)$. \\Do đó, tứ phân vị thứ nhất của mẫu số liệu ghép nhóm là $$Q_1=250+\dfrac{\dfrac{150}{4}-24}{62}\cdot (300-250)=\dfrac{16\,175}{62}.$$
\itemch Sai. Tứ phân vị thứ ba $Q_3$ là $x_{113}\in [300;350)$.\\
Do đó, tứ phân vị thứ ba của mẫu số liệu ghép nhóm là $$Q_3=300+\dfrac{\dfrac{3\cdot150}{4}-(24+62)}{34}\cdot (350-300)=\dfrac{11\,525}{34}.$$
\itemch Sai Khoảng tứ phân vị của mẫu số liệu ghép nhóm là $$\Delta_Q = Q_3-Q_1=\dfrac{11\,525}{34} - \dfrac{16\,175}{62} = \dfrac{41\,150}{527}.$$
\itemch Đúng. Vì doanh nghiệp cần hướng đến các gia đình có mức thu nhập trong khoảng $[Q_1;Q_3)$.
\end{itemchoice}
}
\end{ex}
\begin{ex}%[2D3V1-4]
Giả sử kết quả khảo sát hai khu vực A và B về độ tuổi kết hôn của một số phụ nữ vừa lập gia đình được cho ở bảng sau:
\begin{center}
\begin{tabular}{|c|c|c|c|c|c|}
\hline
\textbf{Tuổi kết hôn} & $[19;22)$ & $[22;25)$ & $[25;28)$ & $[28;31)$ & $[31;34)$\\
\hline
\textbf{Số phụ nữ khu vực A} & $10$ & $27$ & $31$ & $25$ & $7$ \\
\hline
\textbf{Số phụ nữ khu vực B} & $47$ & $40$ & $11$ & $2$ & $0$ \\
\hline
\end{tabular}
\end{center}
\choiceTF
{Khoảng biến thiên của mẫu số liệu ghép nhóm ứng với khu vực A và khu vực B là $15$ tuổi}
{\True Khoảng tứ phân vị của mẫu số liệu ghép nhóm ứng với khu vực A là $\dfrac{388}{75}$}
{\True Khoảng tứ phân vị của mẫu số liệu ghép nhóm ứng với khu vực A là $\dfrac{1\,647}{470}$}
{Phụ nữ ở khu vực A có độ tuổi kết hôn đồng đều hơn}
\loigiai{
\begin{itemchoice}
\itemch Sai. Khoảng biến thiên của mẫu số liệu ghép nhóm ứng với khu vực A là $R=34-19=15$.
\itemch Đúng. Cỡ mẫu $n=10+27+31+25+7=100$.\\
Gọi $x_1$; $x_2$; $\ldots$; $x_{100}$ là mẫu số liệu gốc về độ tuổi kết hôn của một số phụ nữ vừa lập gia đình ở khu vực A được xếp theo thứ tự không giảm.\\
Ta có $x_1$; $x_2$; $\ldots$; $x_{10} \in [19; 22)$, $x_{11}$; $x_{12}$; $\ldots$; $x_{37} \in [22; 25)$, $x_{38}$; $x_{39}$; $\ldots$; $x_{68} \in [25; 28)$, $x_{69}$; $\ldots$; $x_{93} \in [28; 31)$, $x_{94}$; $\ldots$; $x_{100} \in [31; 34)$.\\
Tứ phân vị thứ nhất của mẫu số liệu gốc là $\dfrac{1}{2}(x_{25}+x_{26}) \in [22; 25)$.\\
Do đó, tứ phân thứ nhất của mẫu số liệu ghép nhóm là $$Q_1=22+\dfrac{\frac{100}{4}-10}{27}\cdot (25-22)=\dfrac{71}{3}.$$
Tứ phân vị thứ ba của mẫu số liệu gốc là $\dfrac{1}{2}(x_{75}+x_{76}) \in [28; 31)$.\\
Do đó, tứ phân thứ ba của mẫu số liệu ghép nhóm là $$Q_3=28+\dfrac{\dfrac{3\cdot 100}{4}-(10+27+31)}{25}\cdot(31-28)=\dfrac{721}{25}.$$	
Khoảng tứ phân vị của mẫu số liệu ghép nhóm về độ tuổi kết hôn của một số phụ nữ vừa lập gia đình ở khu vực A là $$\Delta_Q=Q_3-Q_1=\dfrac{721}{25}-\dfrac{71}{3}=\dfrac{388}{75}.$$
\itemch Đúng. Cỡ mẫu $n'=47+40+11+2=100$.\\
Gọi $y_1$; $y_2$; $\ldots$; $y_{100}$ là mẫu số liệu gốc về độ tuổi kết hôn của một số phụ nữ vừa lập gia đình ở khu vực B được xếp theo thứ tự không giảm.\\
Ta có $y_1$; $y_2$; $\ldots$; $y_{47} \in [19; 22)$, $y_{48}$; $y_{49}$; $\ldots$; $y_{87} \in [22; 25)$, $y_{88}$; $y_{89}$; $\ldots$; $y_{98} \in [25; 28)$, $y_{99}$; $y_{100} \in [28; 31)$.\\
Tứ phân vị thứ nhất của mẫu số liệu gốc là $\dfrac{1}{1}(y_{25}+y_{26}) \in [19; 22)$.\\
Do đó, tứ phân thứ nhất của mẫu số liệu ghép nhóm là $$Q'_1 = 19 + \dfrac{\dfrac{100}{4}}{47}\cdot (22-19)=\dfrac{968}{47}.$$
Tứ phân vị thứ ba của mẫu số liệu gốc là $\dfrac{1}{2}(y_{75}+y_{76})\in[22; 25)$.\\
Do đó, tứ phân thứ ba của mẫu số liệu ghép nhóm là $$Q'_3=22+\dfrac{\dfrac{3\cdot100}{4}-47}{40}\cdot (25-22)=\dfrac{968}{47}.$$
Khoảng tứ phân vị của mẫu số liệu ghép nhóm về độ tuổi kết hôn của một số phụ nữ vừa lập gia đình ở khu vực B là $$\Delta'_Q = Q'_3-Q'_1=\dfrac{241}{10}-\dfrac{968}{47}=\dfrac{1\,647}{470}.$$
\itemch Sai. Vì $\Delta'_Q < \Delta_Q$ nên phụ nữ ở khu vực B có độ tuổi kết hôn đồng đều hơn.
\end{itemchoice}
}
\end{ex}
\Closesolutionfile{ans}
\indapan{3}{ans/ans-0-B15-DS}
\Opensolutionfile{ans}[ans/ans-0-B15-KQ]
\caukq
\begin{ex}%[2D3H1-2]
	Bảng $8$ biểu diễn mẫu số liệu ghép nhóm về số tiền (đơn vị: nghìn đồng) mà $60$ khách hàng mua sách ở một cửa hàng trong một ngày.
	\begin{center}
		\begin{tabular}{|c|c|c|c|c|c|}
	\hline Nhóm  &$[40 ; 50)$ & {$[50 ; 60)$} & {$[60 ; 70)$} &	{$[70 ; 80)$} &	{$[80 ; 90)$} \\
	\hline Tần số & $3$ & $6$ & $19$ & $23$ &$9$\\
	\hline
\end{tabular}
	\end{center}
Khoảng biến thiên của mẫu số liệu ghép nhóm trên bằng bao nhiêu ?
\shortans{$50$}
\loigiai{Trong mẫu số liệu ghép nhóm ở bảng $8$, ta có đầu mút trái của nhóm $1$ là $a_1=40$, đầu mút phải của nhóm $5$ là $a_6=90$.\\Vậy khoảng biến thiên của mẫu số liệu ghép nhóm đó là
		$R=a_6-a_1=90-40=50.$}
\end{ex}
\begin{ex}%[2D3V1-3]
	Bảng sau biểu diễn mẫu số liệu ghép nhóm về số tiền (đơn vị: nghìn đồng) mà $60$ khách hàng mua sách ở một cửa hàng trong một ngày.
	\begin{center}
		\begin{tabular}{|c|c|c|c|c|c|}
			\hline Nhóm  &$[40 ; 50)$ & {$[50 ; 60)$} & {$[60 ; 70)$} &	{$[70 ; 80)$} &	{$[80 ; 90)$} \\
			\hline Tần số & $3$ & $6$ & $19$ & $23$ &$9$\\
			\hline
		\end{tabular}
	\end{center}
	Khoảng tứ phân vị của mẫu số liệu ghép nhóm trên bằng bao nhiêu (làm tròn đến 1 chữ số thập phân)
	\shortans{$14{,}2$}
	\loigiai{
Ta có số phần tử của mẫu là $n=60$.\\
Áp dụng công thức, ta có tứ phân vị thứ nhất là $$Q_1=60+\left(\dfrac{15-9}{19}\right)\cdot 10=\dfrac{1\,200}{19}~\text{(nghìn đồng)}.$$
Tứ phân vị thứ ba là
		$$Q_3=70+\left(\dfrac{45-28}{23}\right)\cdot10=\dfrac{1\,780}{23} ~\text{(nghìn đồng)}.$$
	Vậy khoảng tứ phân vị của mẫu số liệu ghép nhóm đã cho là 
	$$\Delta _Q=Q_3-Q_1=\dfrac{1\,780}{23}-\dfrac{1\,200}{19}\approx 14{,}2 ~\text{(nghìn đồng)}.$$}
\end{ex}
\begin{ex}%[2D3V1-3]
Điểm thi đánh giá năng lực của một trường Đại học A thể hiển ở biểu đồ cột như sau.
\begin{center}
	\begin{tikzpicture}[scale=1, font=\footnotesize, line join=round, line cap=round, >=stealth]
		\coordinate [label= left: $0$](O) at (0,0) ;
		\coordinate [label= left: Số học sinh](y) at (0,8) ;
		\coordinate [label= below left: Điểm](x) at (14.3,0) ;
		\foreach \i in {2,3,4,5,6,7}
		{\draw[-,color=gray!40] (0,\i)--(14,\i); }
		\draw (0,1)node[left]{$20$}(0,2)node[left]{$40$}(0,3)node[left]{$60$}(0,4)node[left]{$80$}(0,5)node[left]{$100$}(0,6)node[left]{$120$}(0,7)node[left]{$140$};
		%--------------------------------------------
		\foreach \i/\j/\a in 
		{0.25/0/0,0.75/0/0,1.25/0/0,1.75/0/0,2.25/0/0,2.75/0/0,
			3.25/0/0,3.75/0.05/1,4.25/0.4/8,4.75/1.2/24,5.25/2.7/54,
			5.75/4.75/95,6.25/4.75/95,6.75/6.65/133,7.25/6.1/122,	7.75/5.2/104,8.25/3.1/62,8.75/2.75/55,9.25/1.05/21,9.75/0.6/12,
			10.25/0/0,10.75/0/0,11.25/0/0,11.75/0/0}{
			\draw[] ($(\i,\j)+(0,0.2)$) node[scale=0.7]{$\a$};
		}
		%-----------------------------------------------
		\foreach \i/\j in 
		{0/0,0.5/0,1/0,1.5/0,2/0,2.5/0,3/0,
			3.5/0.05,4/0.4,4.5/1.2,5/2.7,5.5/4.75,6/4.75,
			6.5/6.65,7/6.1,7.5/5.2,8/3.1,8.5/2.75,9/1.05,
			9.5/0.6,10/0,10.5/0,11/0,11.5/0}{
			\fill[blue!20] (\i,0) rectangle+ (0.5,\j);
			\draw[] (\i,0) rectangle+ (0.5,\j);
		}
		\foreach \i/\j in 
		{0.25/[0-50],0.75/[51-100],1.25/[101-150],1.75/[151-200],2.25/[201-250],2.75/[251-300],3.25/[301-350],3.75/[351-400],4.25/[401-450],4.75/[451-500],5.25/[501-550],5.75/[551-600],6.25/[601-650],6.75/[651-700],7.25/[701-750],7.75/[751-800],8.25/[801-850],8.75/[851-900],9.25/[901-950],9.75/[951-1000],10.25/[1001-1050],10.75/[1051-1100],11.25/[1101-1150],11.75/[1151-1200]}{
			\draw($(\i,0)+(0,-1)$) node[rotate=90]{\j};
		}	
		\foreach \i/\j in {}{\draw (\i,\j) node[above]{\j};
		}			
		\draw[->] (O)--(x);
		\draw[->] (O)--(y);
	\end{tikzpicture}
\end{center}
Khoảng tứ phân vị của mẫu điểm thi đánh giá năng lực này bằng bao nhiêu (làm tròn đến hàng đơn vị)
\shortans{$170$}
\loigiai{
Ta thống kê theo bảng tần số như sau
	\begin{center}
	\begin{tabular}{|c|c|c|c|c|c|c|c|c|}
		\hline Nhóm  &$[351 ; 400]$ & $[401 ; 450]$ & $[451 ; 500]$ &	$[501 ; 550]$ &	$[551 ; 600])$ & $[601 ; 650]$ & $[651 ; 700]$ & $[701 ; 750]$\\
		\hline Tần số & $1$ & $8$ & $24$ & $54$ &$95$& $95$ & $132$ & $122$ \\
		\hline
	\end{tabular}
\end{center}
	\begin{center}
	\begin{tabular}{|c|c|c|c|c|c|}
		\hline Nhóm  &	$[751 ; 800]$ &	$[801 ; 850]$& $[851 ; 900]$ & $[901 ; 950]$ &	$[951 ; 1000]$ \\
		\hline Tần số & $104$ &$62$& $55$ & $21$ & $12$ \\
		\hline
	\end{tabular}
\end{center}
Cỡ mẫu của bảng số liệu là $n=785$ học sinh.\\
Tứ phân vị thứ nhất là $Q_1=601+\dfrac{\dfrac{785}{4}-182}{95}\times 49=608{,}35$\\
Tứ phân vị thứ ba là $Q_3=751+\dfrac{\dfrac{3}{4}\times 785-531}{104}\times 49=778{,}21$\\
Do đó, khoảng phân vị là $\Delta Q=169{,}86\approx 170$.

}
\end{ex}
\begin{ex}%[2D3H1-2]
Cho bảng biểu diễn mẫu số liệu ghép nhóm thống kê mức lương của một công ty (đơn vị: triệu đồng).
\begin{center}\begin{tabular}{|c|c|c|c|c|c|c|}
			\hline Nhóm & $[10 ; 15)$ & $[15 ; 20)$ & $[20 ; 25)$ & $[25 ; 30)$ & $[30 ; 35)$ & $[35 ; 40)$  \\
			\hline Tần số & $15$ & $18$ & $10$ &$10$ &$5$ &$2$\\
			\hline
\end{tabular}\end{center}
Khoảng biến thiên của mẫu số liệu ghép nhóm đó là
\shortans{$30$}
	\loigiai{
Trong mẫu số liệu ghép nhóm ở bảng $9$, ta có đầu mút trái của nhóm $1$ là $a_1=10$, đầu mút phải của nhóm $6$ là $a_7=40$.\\Vậy khoảng biến thiên của mẫu số liệu ghép nhóm đó là $R=a_7-a_1=40-10=30.$

}
\end{ex}
\begin{ex}%[2D3V1-3]
	Cho bảng biểu diễn mẫu số liệu ghép nhóm thống kê mức lương của một công ty (đơn vị: triệu đồng).
\begin{center}\begin{tabular}{|c|c|c|c|c|c|c|}
		\hline Nhóm & $[10 ; 15)$ & $[15 ; 20)$ & $[20 ; 25)$ & $[25 ; 30)$ & $[30 ; 35)$ & $[35 ; 40)$  \\
		\hline Tần số & $15$ & $18$ & $10$ &$10$ &$5$ &$2$\\
		\hline
\end{tabular}\end{center}
Khoảng tứ phân vị của mẫu số liệu ghép nhóm đã cho là 
\shortans{$11$}
\loigiai{Số phần tử của mẫu là  $n=60$.\\
Ta có $\dfrac{n}{4}=\dfrac{60}{4}=15$ mà $15<33$.\\
Áp dụng công thức, ta có tứ phân vị thứ nhất là $$Q_1=15+\left(\dfrac{15-15}{18}\right)\cdot 5=15 ~\text{(triệu đồng)}.$$
Ta có $\dfrac{3n}{4}=\dfrac{3\cdot 60}{4}=45$ mà $43<45<53$. Áp dụng công thức, ta có tứ phân vị thứ ba là
	$$Q_3=25+\left(\dfrac{45-43}{10}\right)\cdot5=26 ~\text{(triệu đồng)}.$$
			Vậy khoảng tứ phân vị của mẫu số liệu ghép nhóm đã cho là 
			$$\Delta _Q=Q_3-Q_1=26-15=11 ~\text{(triệu đồng)}.$$
			

	}
\end{ex}
\begin{ex}%[2D3H1-2]
Cho bảng biểu diễn mẫu số liệu ghép nhóm về độ tuổi của cư dân trong một khu phố.
	\begin{center}	\begin{tabular}{|c|c|c|c|c|c|c|c|}
			\hline Nhóm & $[20 ; 30)$& $[30 ; 40)$ & $[40 ; 50)$ & $[50 ; 60)$& $[60 ; 70)$ & $[70 ; 80)$ \\ 
			\hline Tần số & $25$ &$20$&$20$ &$15$&$14$&$6$ \\
			\hline
	\end{tabular}\end{center}
Tính khoảng biến thiên của mẫu số liệu ghép nhóm đó.
\shortans{$60$}
\loigiai{
Trong mẫu số liệu ghép nhóm ở bảng $10$, ta có đầu mút trái của nhóm $1$ là $a_1=20$, đầu mút phải của nhóm $6$ là $a_7=80$.\\Vậy khoảng biến thiên của mẫu số liệu ghép nhóm đó là $R=a_7-a_1=80-20=60$.}
\end{ex}
\Closesolutionfile{ans}
\indapan{6}{ans/ans-0-B15-KQ}